```latex
\documentclass[10pt,twocolumn]{article}

% ------------------------------------------------------------------
% IEEE/top-conference-inspired compact report template
% For COMP2050 AI Programming Project / PaperOrchestra workflow.
%
% PaperOrchestra/Codex instructions:
% - Fill TODO placeholders using workspace/inputs/idea.md,
%   workspace/inputs/experimental_log.md, and conference_guidelines.md.
% - Do not invent metrics, datasets, citations, or experimental results.
% - All numeric claims must come from the metrics files listed in
%   experimental_log.md.
% ------------------------------------------------------------------

% Encoding and typography
\usepackage[utf8]{inputenc}
\usepackage[T1]{fontenc}
\usepackage{times}
\usepackage{microtype}

% Page layout: compact two-column style
\usepackage[
    letterpaper,
    margin=0.72in,
    columnsep=0.24in
]{geometry}

% Math and symbols
\usepackage{amsmath}
\usepackage{amssymb}
\usepackage{mathtools}

% Figures, tables, and captions
\usepackage{graphicx}
\usepackage{booktabs}
\usepackage{multirow}
\usepackage{array}
\usepackage{caption}
\usepackage{subcaption}

% References and links
\usepackage[numbers,sort&compress]{natbib}
\usepackage[hidelinks]{hyperref}
\usepackage{xcolor}

% Lists and spacing
\usepackage{enumitem}
\setlist{nosep,leftmargin=*}

% Section spacing: compact conference-like layout
\usepackage{titlesec}
\titlespacing*{\section}{0pt}{0.85em}{0.45em}
\titlespacing*{\subsection}{0pt}{0.65em}{0.35em}
\titlespacing*{\subsubsection}{0pt}{0.5em}{0.25em}

% Caption style
\captionsetup{
    font=small,
    labelfont=bf,
    skip=4pt
}

% Reduce float spacing
\setlength{\textfloatsep}{8pt plus 1pt minus 2pt}
\setlength{\floatsep}{6pt plus 1pt minus 2pt}
\setlength{\intextsep}{6pt plus 1pt minus 2pt}

% Optional helpers
\newcommand{\methodname}{Budget-Aware Early Exiting}
\newcommand{\dataset}{CIFAR-10}
\newcommand{\backbone}{ResNet-18}

\title{\vspace{-1.0em}
Learning When to Stop: Budget-Aware Early Exiting for Energy-Efficient CNN Inference
\vspace{-0.4em}
}

\author{
Tran Anh Chuong\\
VinUniversity\\
COMP2050 -- Artificial Intelligence Programming Project\\
\texttt{TODO: student email if required}
}

\date{}

\begin{document}

\twocolumn[
\maketitle

\begin{abstract}
TODO(section-writing-agent): Write a one-paragraph abstract summarizing the problem, the proposed early-exit ResNet-18 approach, the compared exit policies, the CIFAR-10 experimental setup, and the main accuracy-efficiency finding. Use only verified metrics from experimental_log.md.
\end{abstract}

\vspace{0.8em}
]

% ------------------------------------------------------------------
\section{Introduction}
% ------------------------------------------------------------------

TODO(literature-review-agent): Motivate Green AI, efficient inference, adaptive computation, and early exiting. Explain why full CNN inference can waste computation on easy inputs. Introduce the central question: whether a CNN can learn when to stop computation to reduce FLOPs, latency, and estimated emissions while preserving accuracy.

TODO(section-writing-agent): State the project contributions clearly. Suggested contribution framing:
\begin{itemize}
    \item an early-exit \backbone{} architecture for \dataset{};
    \item a comparison of fixed-threshold, dynamic-threshold, supervised-controller, and REINFORCE-based exit policies;
    \item an empirical evaluation of accuracy, FLOPs, latency, estimated CO$_2$, and exit behavior;
    \item an analysis of the accuracy-efficiency tradeoff and limitations.
\end{itemize}

% ------------------------------------------------------------------
\section{Related Work}
% ------------------------------------------------------------------

TODO(literature-review-agent): Cover efficient neural inference, early-exit networks, dynamic neural networks, adaptive computation, confidence-based exiting, learned routing/controller policies, and energy-aware or Green AI. Use verified citations only. Do not invent references.

% ------------------------------------------------------------------
\section{Methodology}
% ------------------------------------------------------------------

TODO(section-writing-agent): Explain the full pipeline and algorithmic variations. Use subsections below.

\subsection{Baseline Model}

TODO(section-writing-agent): Describe the full \backbone{} classifier on \dataset{}. Explain that this baseline performs a complete forward pass for every image and serves as the reference for accuracy, FLOPs, latency, and emissions.

\subsection{Early-Exit Architecture}

TODO(section-writing-agent): Describe the modified \backbone{} with auxiliary classifiers after intermediate stages. Explain that each exit outputs logits, probabilities, confidence, and a prediction.

For exit $i$, the probability vector is:
\begin{equation}
    p_i = \operatorname{softmax}(z_i),
\end{equation}
where $z_i$ is the logit vector at exit $i$.

The confidence score is:
\begin{equation}
    c_i = \max_j p_i(j).
\end{equation}

TODO(section-writing-agent): Explain the joint weighted loss across exits:
\begin{equation}
    \mathcal{L}
    =
    \sum_{i=1}^{K} \alpha_i \mathcal{L}_i,
\end{equation}
where $K$ is the number of exits and $\alpha_i$ controls the importance of each exit loss.

\subsection{Fixed Confidence Thresholding}

TODO(section-writing-agent): Explain the fixed-threshold exit rule. Exit at the first intermediate classifier whose confidence exceeds threshold $\tau$:
\begin{equation}
    \text{exit at } i
    \quad \text{if} \quad
    c_i \geq \tau.
\end{equation}

TODO(section-writing-agent): Mention the evaluated threshold values only if verified in experimental_log.md.

\subsection{Budget-Aware Dynamic Thresholding}

TODO(section-writing-agent): Explain the dynamic threshold policy. Define the normalized FLOPs consumed up to exit $i$:
\begin{equation}
    r_i =
    \frac{\operatorname{FLOPs}_i}
    {\operatorname{FLOPs}_{\text{full}}}.
\end{equation}

TODO(section-writing-agent): Describe the accuracy-first and budget-first variants if they were implemented and evaluated. Use verified equations from idea.md or experimental_log.md.

\subsection{Supervised Exit Controller}

TODO(section-writing-agent): Explain the learned controller. At each exit, the controller observes a state vector:
\begin{equation}
    s_i = [c_i, H_i, m_i, d_i, r_i],
\end{equation}
where $c_i$ is confidence, $H_i$ is entropy, $m_i$ is top-1/top-2 margin, $d_i$ is normalized exit depth, and $r_i$ is normalized FLOPs.

The predictive entropy is:
\begin{equation}
    H_i = -\sum_j p_i(j)\log p_i(j).
\end{equation}

The margin is:
\begin{equation}
    m_i = p_i(\text{top-1}) - p_i(\text{top-2}).
\end{equation}

TODO(section-writing-agent): Explain whether the controller was trained using earliest-correct-exit labels, reward labels, or another verified supervision signal.

\subsection{REINFORCE Exit Controller}

TODO(section-writing-agent): Explain the policy-gradient controller if included in the final experiments. The controller policy is:
\begin{equation}
    \pi_\theta(a_i \mid s_i),
\end{equation}
where $a_i \in \{\text{exit}, \text{continue}\}$.

Use the reward:
\begin{equation}
    R =
    \mathbb{1}[\hat{y}=y]
    -
    \lambda
    \frac{\operatorname{FLOPs}_i}
    {\operatorname{FLOPs}_{\text{full}}}.
\end{equation}

TODO(section-writing-agent): Explain that the early-exit backbone is frozen during REINFORCE training if this is verified by the implementation.

% ------------------------------------------------------------------
\section{Experimental Setup}
% ------------------------------------------------------------------

TODO(section-writing-agent): Describe the dataset, training setup, evaluation protocol, and metrics. Use only information from idea.md, experimental_log.md, configs, and run manifests.

\subsection{Dataset}

TODO(section-writing-agent): Describe \dataset{} as the image classification benchmark used in the project. Mention train/test split only if verified.

\subsection{Compared Methods}

TODO(section-writing-agent): List the compared methods:
\begin{itemize}
    \item Full \backbone{} baseline;
    \item Early-exit \backbone{};
    \item fixed-threshold policy;
    \item dynamic-threshold policy;
    \item supervised learned controller;
    \item REINFORCE controller, if evaluated.
\end{itemize}

\subsection{Evaluation Metrics}

TODO(section-writing-agent): Define the metrics used in the report: accuracy, FLOPs/sample, FLOPs reduction, latency/sample, estimated CO$_2$ emissions, exit distribution, average exit, and reward if applicable.

FLOPs reduction is computed as:
\begin{equation}
    \operatorname{Reduction}
    =
    1 -
    \frac{\operatorname{FLOPs}_{\text{method}}}
    {\operatorname{FLOPs}_{\text{full}}}.
\end{equation}

\subsection{Statistical Analysis}

TODO(section-writing-agent): State whether multiple runs, repeated seeds, or robustness analysis were performed. If not, justify the limitation according to conference_guidelines.md. Do not claim statistical significance unless supported by the metrics.

% ------------------------------------------------------------------
\section{Results}
% ------------------------------------------------------------------

TODO(section-writing-agent): Insert the main comparison table using the source-of-truth metrics from experimental_log.md. Do not invent numbers.

\begin{table*}[t]
\centering
\caption{Main comparison of full and adaptive inference methods. TODO(section-writing-agent): Fill values from the source-of-truth metrics file only.}
\label{tab:main-results}
\begin{tabular}{lrrrrrl}
\toprule
Method & Accuracy & FLOPs/sample & FLOPs Red. & Latency/sample & CO$_2$ & Avg. Exit \\
\midrule
Full \backbone{} & TODO & TODO & 0\% & TODO & TODO & Final \\
Early-Exit \backbone{} & TODO & TODO & TODO & TODO & TODO & TODO \\
Fixed Threshold & TODO & TODO & TODO & TODO & TODO & TODO \\
Dynamic Threshold & TODO & TODO & TODO & TODO & TODO & TODO \\
Supervised Controller & TODO & TODO & TODO & TODO & TODO & TODO \\
REINFORCE Controller & TODO & TODO & TODO & TODO & TODO & TODO \\
\bottomrule
\end{tabular}
\end{table*}

TODO(plotting-agent): Insert or reference the final accuracy comparison figure if available.

\begin{figure}[t]
    \centering
    \includegraphics[width=\linewidth]{figures/accuracy_comparison.png}
    \caption{TODO(plotting-agent): Caption based on source metrics. Do not claim improvement unless supported.}
    \label{fig:accuracy-comparison}
\end{figure}

TODO(plotting-agent): Insert or reference the final accuracy versus FLOPs figure if available.

\begin{figure}[t]
    \centering
    \includegraphics[width=\linewidth]{figures/accuracy_vs_flops.png}
    \caption{TODO(plotting-agent): Caption describing the accuracy--FLOPs tradeoff using verified data.}
    \label{fig:accuracy-flops}
\end{figure}

TODO(plotting-agent): Insert or reference the final accuracy versus latency figure if available.

\begin{figure}[t]
    \centering
    \includegraphics[width=\linewidth]{figures/accuracy_vs_latency.png}
    \caption{TODO(plotting-agent): Caption describing the accuracy--latency tradeoff using verified data.}
    \label{fig:accuracy-latency}
\end{figure}

TODO(plotting-agent): Insert or reference the CO2 or energy-efficiency figure if available.

\begin{figure}[t]
    \centering
    \includegraphics[width=\linewidth]{figures/co2_reduction.png}
    \caption{TODO(plotting-agent): Caption describing estimated CO$_2$ or energy-efficiency results. State if values are estimates.}
    \label{fig:co2-reduction}
\end{figure}

\subsection{Ablation Analysis}

TODO(section-writing-agent): Discuss threshold ablations, dynamic threshold ablations, controller lambda sweep, or REINFORCE reward sweep if supported by the metrics.

\begin{table}[t]
\centering
\caption{TODO(section-writing-agent): Ablation study summary. Fill only with verified values.}
\label{tab:ablation}
\begin{tabular}{lrrr}
\toprule
Setting & Accuracy & FLOPs Red. & Notes \\
\midrule
TODO & TODO & TODO & TODO \\
TODO & TODO & TODO & TODO \\
TODO & TODO & TODO & TODO \\
\bottomrule
\end{tabular}
\end{table}

% ------------------------------------------------------------------
\section{Discussion}
% ------------------------------------------------------------------

TODO(section-writing-agent): Interpret the results. Discuss the accuracy-efficiency tradeoff, why the best fixed-threshold policy may work well, why the supervised controller may exit too aggressively, and why REINFORCE improves over the supervised controller if supported by the results.

TODO(section-writing-agent): Clearly separate verified conclusions from hypotheses. Do not overclaim beyond CIFAR-10 or beyond the measured hardware/metrics.

% ------------------------------------------------------------------
\section{Limitations}
% ------------------------------------------------------------------

TODO(section-writing-agent): Discuss limitations:
\begin{itemize}
    \item evaluation limited to \dataset{};
    \item latency and emissions may depend on hardware and measurement setup;
    \item estimated CO$_2$ values should be interpreted cautiously unless directly measured;
    \item repeated-seed or variance analysis may be limited;
    \item adaptive inference adds policy/controller complexity.
\end{itemize}

% ------------------------------------------------------------------
\section{Conclusion}
% ------------------------------------------------------------------

TODO(section-writing-agent): Summarize the project, the early-exit approach, the compared policies, the main empirical findings, and future work. Mention that negative or mixed results are still valuable if they reveal the tradeoffs between accuracy and efficiency.

% ------------------------------------------------------------------
\section*{Acknowledgements}
% ------------------------------------------------------------------

TODO(section-writing-agent): Acknowledge any allowed help, debugging support, tools, or state that no external help was received. Mention AI assistance only according to the project statement and AI usage policy.

% ------------------------------------------------------------------
% References
% ------------------------------------------------------------------

\bibliographystyle{plainnat}
\bibliography{refs}

\end{document}
```
